% abstract_intro.tex  -- written by Harshvardhan
\begin{abstract}
Matrix multiplication dominates the arithmetic of modern machine-learning
workloads, and general-purpose processors execute it inefficiently because
every multiply-accumulate (MAC) operation re-fetches its operands. Systolic
arrays, the structure at the core of Google's Tensor Processing Unit, reuse
each operand across a row or column of processing elements (PEs). This paper
presents the complete design and evaluation of a synthesizable, parameterised
$N\times N$ output-stationary systolic-array accelerator in Verilog-2005,
built around signed Q4.4 operands and 32-bit accumulators. Phase~I delivered
a verified $4\times4$ prototype. Phase~II adds row-wide double-buffered
(ping-pong) operand buffers, a result buffer, back-to-back and accumulate
modes, output requantisation with saturation, hardware performance
counters, optional pipeline registers, a sequential single-MAC baseline
and a weight-stationary comparison array. More than 85\,000 element-level
checks, including a 1000-matrix random regression, pass against a bit-exact
Python golden model. A row-wide interface and overlapping data transfer
with computation bring the cost of a $4\times4$ product from 60 to 12
cycles. Two pipeline registers
raise the post-route clock on a Lattice ECP5 FPGA from 51 to 116~MHz, so a
$4\times4$ product completes every 121~ns: 12.4$\times$ faster than the
sequential baseline, using 16 DSP blocks. A cycle-exact analytical model
predicts every measured latency, and the comparison with a
weight-stationary array explains why output-stationary utilisation is
capped at one third.
\end{abstract}

\begin{IEEEkeywords}
systolic array, matrix multiplication, processing element, MAC,
fixed-point arithmetic, double buffering, FPGA, Verilog
\end{IEEEkeywords}

% -----------------------------------------------------------------------------
\section{Introduction}
As Dennard scaling ended, performance growth moved from faster general-purpose
cores to \emph{domain-specific architectures}~\cite{hennessy2019golden}. Deep
learning is the most prominent target. Its fully connected and convolutional
layers reduce to general matrix multiplication (GEMM)~\cite{sze2017efficient}.
On a CPU, a GEMM triple loop is limited by the memory wall: each MAC needs two
operand loads. A \emph{systolic array}~\cite{kung1979systolic,kung1982why}
instead passes operands between neighbouring PEs, so an operand fetched once
feeds a whole row or column. The result is high throughput, short wires and
a regular structure.

This project designs, verifies and evaluates a small systolic GEMM
accelerator, in line with Module~6 of the course (heterogeneous computing,
TPU/AI hardware). It builds on the digital-design foundations of
Modules~1--2 and on the performance concepts of Modules~4--5 (CPI, Amdahl's
law, pipelining)~\cite{patterson_cod,harris_ddca}. The work is entirely
simulation and synthesis based; no board is required. The contributions are:
\begin{itemize}
\item a parameterised output-stationary (OS) array with a complete host
interface, double buffering, accumulate mode and on-chip performance counters;
\item a sequential single-MAC baseline and a weight-stationary (WS) streaming
array, both verified, for quantitative comparison;
\item a verification environment of 10 self-checking testbenches driven by a
bit-exact golden model;
\item post-route timing and resource results, a pipelining study and a
cycle-exact analytical performance model.
\end{itemize}
